\section{Automatic routing and supervision}
\label{sec:routing}

The router decides, for each ontology, which of the procedures of
Section~\ref{sec:procedures} runs, with which options, and which procedure
runs if the first declines.  It reads only the source profile computed by the
frontend: counts of axiom and constructor kinds, distinct names, sizes,
depths, and bounds, the Konclude-compatible expressivity code, and
source-level certificates.  The router reads no ontology identifier and no
expected answer, and \texttt{km profile} prints the profile and the selected
route for a source.

The route is selected after the first parse pass and before normalisation,
because options such as triggered absorption and cardinality metadata change
the normalised clause set.  The decision has three parts.  A fragment gate
first fixes which procedures are semantically applicable.  Every ontology with
DL-safe rules selects the rule route.  Every ontology with an ABox or with
nominals selects a procedure that represents individuals, unless a
source-level certificate shows that the ABox cannot change a named-class
subsumption.  That certificate holds when the source has only positive
assertions and contains no bottom class or role, complement, disjointness,
number restriction, universal role, rule, key, or datatype constraint; for
such a source, interpreting every class and role as full and every individual
as one element is a model, and nominal-free $\mathcal{SRIQ}$ is preserved
under disjoint unions, so the TBox alone determines the taxonomy.  Within the
set of applicable procedures, the policy then selects a procedure from
fragment membership, from the presence of constructs, and from further
source certificates, for example that no number-restricted role shares a
hierarchy component with an inverse, chained, or transitive role.  The size of
the source enters only into resource decisions such as thread counts and
in-process execution.  An input that the policy cannot assign to a specialist
goes to a portfolio of procedures that are designed to be complete for the
input's fragment.  \todo{Check this paragraph against \texttt{docs/ROUTING.md}
and the routing source at v1.7.0, and add the figure and the predicate table
for the general policy; the earlier manuscript documents the v1.5.3 selector,
which the v1.5.4 policy replaced.}

Each route is a named bundle of options that replaces the ambient settings, so
that a later frontend pass cannot be routed a second time.  An atomic route
runs one worker.  When that worker declines, fails, or exceeds its memory
limit, the supervisor reruns the classification under a fallback bundle: the
portfolio for nominal-free sources and sources with a certified ABox, and the
nominal-aware consequence-based engine for every other source.  The portfolio
composes several workers.  A small source that passes the EL clause screen is
classified by EL completion inside the supervisor.  Otherwise the
consequence-based engine runs beside a hypertableau worker and a certified EL
worker in repair mode, and the supervisor accepts the first complete answer
under acceptance rules that prefer the consequence-based answer.  A
finite-cover route handles ontologies that entail $\top\sqsubseteq\{o_1,\ldots,
o_k\}$ for two to four individuals: every model then has at most $k$ elements,
and the route decides the classification by a SAT search over those elements.
With a cover of one individual, the consequence-based engine merges every
individual into the single element.

Every reasoning worker is a child process that reads the clause file and
writes its answer to a temporary file, which the supervisor parses as a
stream, so a taxonomy of several gigabytes is never held in a string or in an
undrained pipe.  A watchdog polls the resident set size of the child and the
wall clock and notices the child's exit through a process descriptor.  A child
that exceeds its memory or time limit is killed and reaped; a portfolio arm
that answers first causes the remaining workers to be killed.  When the
parallel run of the consequence-based engine exceeds its memory limit, the
supervisor retries once on one thread with the per-function successor
schedule.  Transparent huge pages are disabled for the process tree, and the
supervisor returns freed heap pages before it waits for a worker, which
reduces the memory of a small classification by 20 to 30~MiB.

A worker answers in an internal vocabulary.  The supervisor maps every internal
name to its source IRI, discards generated names, drops reflexive pairs, and
collects subsumption by the bottom class into the set of unsatisfiable named
classes.  The result of \texttt{km classify} contains a consistency value, the
sorted named subsumption pairs, the sorted unsatisfiable classes, and the
number of dropped clauses, which is zero for a published answer.  Three
consistency rules combine frontend facts with the taxonomy: an individual
asserted into two disjoint named classes, a clash in the rule stage, or an
unsatisfiable class with an asserted member each makes the ontology
inconsistent.  A worker that reports incompleteness, a decline, or a nonzero
exit status publishes nothing.
